% !TeX program = pdflatex
\documentclass[11pt,a4paper]{amsart}

\newcommand{\notelanguage}{finnish}
% Shared preamble for course notes and exercise sets.

\usepackage[T1]{fontenc}
\usepackage{lmodern}
\providecommand{\notelanguage}{english}
\usepackage[finnish,english]{babel}
\AtBeginDocument{\expandafter\selectlanguage\expandafter{\notelanguage}}

\newcommand{\notesfinnishlanguage}{finnish}
\ifx\notelanguage\notesfinnishlanguage
\newcommand{\notestheoremname}{Lause}
\newcommand{\noteslemmaname}{Lemma}
\newcommand{\notespropositionname}{Propositio}
\newcommand{\notescorollaryname}{Seurauslause}
\newcommand{\notesdefinitionname}{Määritelmä}
\newcommand{\notesexamplename}{Esimerkki}
\newcommand{\notesproblemname}{Tehtävä}
\newcommand{\notesexercisename}{Tehtävä}
\newcommand{\notesremarkname}{Huomautus}
\newcommand{\notessolutionname}{Ratkaisu}
\else
\newcommand{\notestheoremname}{Theorem}
\newcommand{\noteslemmaname}{Lemma}
\newcommand{\notespropositionname}{Proposition}
\newcommand{\notescorollaryname}{Corollary}
\newcommand{\notesdefinitionname}{Definition}
\newcommand{\notesexamplename}{Example}
\newcommand{\notesproblemname}{Problem}
\newcommand{\notesexercisename}{Exercise}
\newcommand{\notesremarkname}{Remark}
\newcommand{\notessolutionname}{Solution}
\fi

\usepackage[a4paper,margin=30mm]{geometry}
\usepackage{microtype}
\usepackage{graphicx}
\usepackage{enumitem}

\newcommand{\exerciseheading}[2]{%
  \par\noindent
  {\large\bfseries #1\par}
  \vspace{0.1em}
  \noindent{\large #2\par}
  \vspace{1.5em}
}

\usepackage{amsmath}
\usepackage{amssymb}
\usepackage{amsthm}
\usepackage{mathtools}
\usepackage{aliascnt}

% Linear systems: the argument gives the number of alignment pairs.
\newenvironment{linearsystem}[1]
  {\left\{\begin{alignedat}{#1}}
  {\end{alignedat}\right.}

\usepackage[hidelinks,pdfusetitle]{hyperref}

\numberwithin{equation}{section}
\setlist{itemsep=0.25em,topsep=0.5em}

\theoremstyle{plain}
\newtheorem{theorem}{\notestheoremname}[section]
\newaliascnt{lemma}{theorem}
\newtheorem{lemma}[lemma]{\noteslemmaname}
\aliascntresetthe{lemma}
\newaliascnt{proposition}{theorem}
\newtheorem{proposition}[proposition]{\notespropositionname}
\aliascntresetthe{proposition}
\newaliascnt{corollary}{theorem}
\newtheorem{corollary}[corollary]{\notescorollaryname}
\aliascntresetthe{corollary}

\theoremstyle{definition}
\newaliascnt{definition}{theorem}
\newtheorem{definition}[definition]{\notesdefinitionname}
\aliascntresetthe{definition}
\newaliascnt{example}{theorem}
\newtheorem{example}[example]{\notesexamplename}
\aliascntresetthe{example}
\newaliascnt{problem}{theorem}
\newtheorem{problem}[problem]{\notesproblemname}
\aliascntresetthe{problem}
\newtheorem{exercise}{\notesexercisename}

\theoremstyle{remark}
\newaliascnt{remark}{theorem}
\newtheorem{remark}[remark]{\notesremarkname}
\aliascntresetthe{remark}

\newenvironment{solution}
{
\begin{proof}[\notessolutionname]}
  {
\end{proof}}

\usepackage[nameinlink,noabbrev]{cleveref}

\crefname{theorem}{theorem}{theorems}
\crefname{lemma}{lemma}{lemmas}
\crefname{proposition}{proposition}{propositions}
\crefname{corollary}{corollary}{corollaries}
\crefname{definition}{definition}{definitions}
\crefname{example}{example}{examples}
\crefname{problem}{problem}{problems}
\crefname{exercise}{exercise}{exercises}
\crefname{remark}{remark}{remarks}

\ifx\notelanguage\notesfinnishlanguage
\crefname{theorem}{lause}{lauseet}
\crefname{lemma}{lemma}{lemmat}
\crefname{proposition}{propositio}{propositiot}
\crefname{corollary}{seurauslause}{seurauslauseet}
\crefname{definition}{määritelmä}{määritelmät}
\crefname{example}{esimerkki}{esimerkit}
\crefname{problem}{tehtävä}{tehtävät}
\crefname{exercise}{tehtävä}{tehtävät}
\crefname{remark}{huomautus}{huomautukset}
\fi

\newcommand{\NN}{\mathbb{N}}
\newcommand{\NNone}{\mathbb{N}_1}
\newcommand{\ZZ}{\mathbb{Z}}
\newcommand{\QQ}{\mathbb{Q}}
\newcommand{\RR}{\mathbb{R}}
\newcommand{\CC}{\mathbb{C}}
\newcommand{\PP}{\mathbb{P}}

\DeclareMathOperator{\im}{im}
\DeclareMathOperator{\rank}{rank}
\DeclareMathOperator{\syt}{syt}
\DeclarePairedDelimiter{\abs}{\lvert}{\rvert}
\DeclarePairedDelimiter{\norm}{\lVert}{\rVert}
\newcommand{\set}[1]{\left\lbrace#1\right\rbrace}
\newcommand{\maxset}[1]{\max\set{#1}}
\newcommand{\minset}[1]{\min\set{#1}}
\newcommand{\powerset}[1]{\mathcal{P}\!\left(#1\right)}


\title{Johdatus lukuteoriaan ja sen sovelluksiin}
\author{Joona Piirainen}
\date{Syksy 2026 -- periodi 2}

\begin{document}

\begin{abstract}
  Luentomuistiinpanot Helsingin yliopiston kurssilta
  Johdatus lukuteoriaan ja sen sovelluksiin (MFK-M213).
  Kurssi järjestetään aikavälillä 27.10.2026--17.12.2026.
\end{abstract}

\maketitle

\tableofcontents

\section{Kurssitiedot}

Kurssin viralliset tiedot löytyvät
\href{https://studies.helsinki.fi/kurssit/opintojakso/otm-989530ba-4d33-4184-bfc4-3d0ed8799643/MFK-M213}{Helsingin
yliopiston kurssisivulta}.

\section{Muistiinpanot}

\subsection{27.10.2026}

\noindent\textit{Aihealue: monisteen luku 2.1 (s.~9--11) ja luku 4
  lauseeseen 39 asti (s.~21--22). Alla mainitut lausenumerot viittaavat
monisteeseen.}

\subsubsection{Jaollisuus}

Kokonaisluku $a$ \emph{jakaa} kokonaisluvun $b$, jos
\[
  a\mid b \quad\Longleftrightarrow\quad
  b=ac \text{ jollakin } c\in\ZZ.
\]
Tällöin $a$ on luvun $b$ tekijä ja $b$ on luvun $a$ monikerta.
Merkintä $a\nmid b$ tarkoittaa, ettei $a$ jaa lukua $b$.
Esimerkiksi $3\mid 24$, koska $24=3\cdot 8$, mutta $5\nmid 24$.

\paragraph{Nollan erikoistapaus.}
Jokaiselle $a\in\ZZ$ pätee $a\mid 0$, myös $0\mid 0$.
Sen sijaan $0\mid b$ vain, kun $b=0$.
Ehto $b/a\in\ZZ$ on siis jaollisuuden kanssa yhtäpitävä vain, kun
$a\ne 0$; määritelmä $b=ac$ toimii myös nollalle.

\paragraph{Jaollisuuden laskusäännöt (lause 5).}
Kun kaikki luvut ovat kokonaislukuja, pätee:
\begin{enumerate}[label=(\alph*)]
  \item Jos $a\mid b$ ja $a\mid c$, niin $a\mid (xb+yc)$ kaikilla
    $x,y\in\ZZ$. Yhteinen tekijä jakaa siis myös summan ja erotuksen.
  \item Jos $a\mid b$ ja $c\mid d$, niin $ac\mid bd$.
  \item Jos $a\mid b$ ja $b\mid c$, niin $a\mid c$.
  \item Jos $a\mid b$ ja $b\mid a$, niin $a=\pm b$.
  \item Jos $a\mid b$ ja $b\ne 0$, niin $|a|\le |b|$.
\end{enumerate}

\begin{proof}
  (a) Kirjoitetaan $b=ak$ ja $c=a\ell$, missä $k,\ell\in\ZZ$.
  Silloin $xb+yc=a(xk+y\ell)$ ja $xk+y\ell\in\ZZ$, joten
  $a\mid(xb+yc)$.

  (b) Oletuksista saadaan $b=ak$ ja $d=c\ell$ joillakin
  $k,\ell\in\ZZ$. Siis $bd=ac(k\ell)$, joten $ac\mid bd$.

  (c) Kirjoitetaan $b=ak$ ja $c=b\ell$, missä $k,\ell\in\ZZ$.
  Tällöin $c=a(k\ell)$, joten $a\mid c$.

  (d) Kirjoitetaan $b=ak$ ja $a=b\ell$, missä $k,\ell\in\ZZ$.
  Jos $b=0$, jälkimmäinen yhtälö antaa $a=0$, joten $a=b$.
  Jos $b\ne0$, sijoittamalla saadaan $b=b\ell k$ ja siten
  $k\ell=1$. Kokonaisluvuille tämä on mahdollista vain, kun
  $k=\ell=1$ tai $k=\ell=-1$. Siis $a=b$ tai $a=-b$.

  (e) Kirjoitetaan $b=ak$, missä $k\in\ZZ$.
  Koska $b\ne0$, myös $k\ne0$, joten $|k|\ge1$.
  Näin ollen $|b|=|a||k|\ge|a|$.
\end{proof}

\subsubsection{Jakoyhtälö ja jakojäännös}

\paragraph{Jakoyhtälö (lause 6).}
Kun $a,b\in\ZZ$ ja $b\ne 0$, on olemassa yksikäsitteiset
$q,r\in\ZZ$, joille
\[
  \boxed{a=qb+r,\qquad 0\le r<|b|.}
\]
Luku $q$ on osamäärä ja $r$ jakojäännös.
Jakojäännös valitaan aina epänegatiiviseksi, vaikka $a$ tai $b$
olisi negatiivinen.

\begin{proof}
  Kun $b>0$, valitse $q=\lfloor a/b\rfloor$ ja aseta $r=a-qb$.
  Ehdosta $qb\le a<(q+1)b$ seuraa $0\le r<b$.
  Kun $b<0$, jo todistetun positiivisen jakajan tapauksen perusteella
  on olemassa $Q,r\in\ZZ$, joille $a=Q|b|+r$ ja $0\le r<|b|$.
  Koska $|b|=-b$, valinta $q=-Q$ antaa $a=qb+r$.

  Jos $a=q_1b+r_1=q_2b+r_2$, niin
  $(q_1-q_2)b=r_2-r_1$. Oikean puolen itseisarvo on pienempi
  kuin $|b|$, joten se voi olla luvun $b$ monikerta vain ollessaan
  nolla. Siis $r_1=r_2$ ja $q_1=q_2$.
\end{proof}

\paragraph{Esimerkkejä.}
\[
  \begin{aligned}
    17 &= 3\cdot 5+2,    & q&=3,  & r&=2,\\
    -17 &= (-4)\cdot 5+3,& q&=-4, & r&=3,\\
    17 &= (-3)\cdot(-5)+2,& q&=-3,& r&=2.
  \end{aligned}
\]
Esitys $-17=(-3)\cdot 5-2$ ei ole vaadittu jakoyhtälö, koska
jakojäännös olisi negatiivinen.

\paragraph{Seurauksia.}
Kun $b\ne 0$, pätee $b\mid a$ täsmälleen silloin, kun $r=0$.
Jakamalla kahdella nähdään, että jokainen kokonaisluku on joko
muotoa $2k$ (parillinen) tai $2k+1$ (pariton), missä $k\in\ZZ$.
Jakajaa $b=0$ ei sallita: ehto $0\le r<|b|$ olisi mahdoton.

\newpage
\subsubsection{Kongruenssit}

Olkoon tässä $n$ positiivinen kokonaisluku.
Kokonaisluvut $a$ ja $b$ ovat \emph{kongruentit modulo $n$}, jos
\[
  \boxed{a\equiv b\pmod n
  \quad\Longleftrightarrow\quad n\mid(a-b).}
\]
Toisin sanoen $a=b+kn$ jollakin $k\in\ZZ$.
Luku $n$ on \emph{modulo}.
Esimerkiksi $3\equiv7\pmod4$ ja $-2\equiv3\pmod5$.

\paragraph{Yhteys jakojäännöksiin.}
Kongruenssi $a\equiv b\pmod n$ tarkoittaa, että luvuilla on sama
jakojäännös jaettaessa luvulla $n$. Nimittäin jos
$a=qn+r$ ja $b=sn+t$, missä $0\le r,t<n$, niin
\[
  n\mid(a-b)\quad\Longleftrightarrow\quad n\mid(r-t)
  \quad\Longleftrightarrow\quad r=t.
\]
Viimeinen ekvivalenssi seuraa ehdosta $|r-t|<n$.
Erityisesti $a\equiv0\pmod n$ tarkoittaa samaa kuin $n\mid a$.

\paragraph{Laskusäännöt (lause 38).}
Jos $a\equiv b\pmod n$ ja $c\equiv d\pmod n$, niin
\[
  a+c\equiv b+d\pmod n,\qquad
  a-c\equiv b-d\pmod n,\qquad
  ac\equiv bd\pmod n.
\]
Laskuissa luvun saa siis korvata sen kanssa kongruentilla,
helpommalla luvulla. Kaikissa kongruensseissa modulo on sama.

\begin{proof}
  Oletusten mukaan $n\mid(a-b)$ ja $n\mid(c-d)$.
  Jaollisuuden laskusäännön (a) perusteella $n$ jakaa luvut
  \[
    \begin{aligned}
      (a+c)-(b+d)&=(a-b)+(c-d),\\
      (a-c)-(b-d)&=(a-b)-(c-d),\\
      ac-bd&=c(a-b)+b(c-d).
    \end{aligned}
  \]
  Kongruenssin määritelmä antaa siis vastaavasti
  $a+c\equiv b+d\pmod n$, $a-c\equiv b-d\pmod n$ ja
  $ac\equiv bd\pmod n$.
\end{proof}

\paragraph{Potenssit (lause 39).}
Jos $a\equiv b\pmod n$, niin kaikilla kokonaisluvuilla $m\ge0$
\[
  a^m\equiv b^m\pmod n.
\]
\begin{proof}
  Todistetaan väite induktiolla eksponentin $m$ suhteen.
  Kun $m=0$, molemmat potenssit ovat $1$, kun eksponentin $0$
  potenssi tulkitaan tyhjäksi tuloksi. Siis $a^0\equiv b^0\pmod n$,
  koska $n\mid(1-1)$.

  Oletetaan, että $a^k\equiv b^k\pmod n$ jollakin $k\ge0$.
  Koska myös $a\equiv b\pmod n$, lauseen 38 kertolaskusääntö antaa
  \[
    a^{k+1}=a^k a\equiv b^k b=b^{k+1}\pmod n.
  \]
  Induktioperiaatteen nojalla väite pätee kaikilla $m\ge0$.
\end{proof}

\paragraph{Laskuesimerkkejä.}
Koska $17\equiv2\pmod5$ ja $14\equiv-1\pmod5$, saadaan
\[
  17\cdot14+8\equiv2\cdot(-1)+3=1\pmod5.
\]
Suuriakin potensseja voidaan käsitellä laskematta niitä auki:
\[
  7^{20}\equiv(-1)^{20}=1\pmod8,
  \qquad
  3^{10}=(3^2)^5\equiv(-1)^5\equiv4\pmod5.
\]

\paragraph{Muista.}
Jaollisuus todistetaan esittämällä luku kokonaislukumonikertana;
kongruenssi todistetaan osoittamalla puolten erotus jaolliseksi
modulolla. Kongruensseissa yhteen-, vähennys- ja kertolasku sekä
kokonaislukupotenssit yllä olevin ehdoin toimivat. Jakaminen ei
yleensä toimi: $2\equiv6\pmod4$, mutta $1\not\equiv3\pmod4$.

\newpage
\subsection{29.10.2026: alkuluvut ja aritmetiikan peruslause}

\subsubsection{Kertaus: suuren potenssin viimeinen numero}
Luvun viimeinen numero on sen jakojäännös modulo $10$.
Koska $3^2=9\equiv-1\pmod{10}$, saadaan
\[
  3^{2023}=3(3^2)^{1011}
  \equiv3(-1)^{1011}=-3\equiv7\pmod{10}.
\]
Vaihtoehtoisesti $3^4\equiv1\pmod{10}$ ja $2023=4\cdot505+3$, joten
$3^{2023}\equiv3^3\equiv7\pmod{10}$.
\emph{Eksponentin saa tässä supistaa modulo $4$, koska $3^4\equiv1$;
yleensä eksponenttia ei saa supistaa samalla modulolla kuin kantalukua.}

\subsubsection{Alkuluvut ja alkutekijät}
\begin{definition}
  Kokonaisluku $p\ge2$ on \emph{alkuluku}, jos sen ainoat
  positiiviset tekijät ovat $1$ ja $p$. Kokonaisluku $n\ge2$ on
  \emph{yhdistetty}, jos se ei ole alkuluku, eli $n=ab$ joillakin
  kokonaisluvuilla $a,b\ge2$. Alkuluku $p$ on luvun $n$
  \emph{alkutekijä}, jos $p\mid n$.
\end{definition}
Esimerkiksi $2,3,5,7,11$ ovat alkulukuja ja $12=3\cdot4$ on
yhdistetty. Luku $1$ ei ole alkuluku eikä yhdistetty luku.
Alkuluvun negatiiviset tekijät ovat $-1$ ja $-p$.

\begin{lemma}[Alkutekijän olemassaolo]
  Jokaisella kokonaisluvulla $n\ge2$ on alkutekijä.
\end{lemma}
\begin{proof}
  Olkoon $d>1$ luvun $n$ pienin positiivinen tekijä; tällainen on
  olemassa, koska $n\mid n$. Jos $d=ab$ joillakin $a,b\ge2$, niin
  $1<a<d$ ja $a\mid n$, mikä on ristiriita minimaalisuuden kanssa.
  Siis $d$ on alkuluku.
\end{proof}

\begin{theorem}[Alkulukujen äärettömyys]
  Alkulukuja on äärettömän monta.
\end{theorem}
\begin{proof}
  Oletetaan, että kaikki alkuluvut ovat $p_1,\ldots,p_k$.
  Luvulla $N=p_1\cdots p_k+1>1$ on alkutekijä $q$.
  Jos $q=p_i$ jollakin $i$, niin $q$ jakaa sekä luvun $N$ että
  tulon $p_1\cdots p_k$ ja siis niiden erotuksen $1$. Tämä on
  mahdotonta, joten $q$ puuttuu luettelosta. Ristiriita.
\end{proof}
Luvun $N$ ei tarvitse itse olla alkuluku. Esimerkiksi
$2\cdot3\cdot5\cdot7\cdot11\cdot13+1=30031=59\cdot509$.
Todistus tarvitsee vain yhden luvun $N$ alkutekijän.

\subsubsection{Yksikäsitteisyyden apuväline: Eukleideen lemma}
Seuraavat apulauseet tekevät aritmetiikan peruslauseen todistuksesta
itsenäisen: yksikäsitteistä alkutekijähajotelmaa ei oleteta etukäteen.

\begin{lemma}[Bézout'n identiteetti]
  Olkoot $a,b\in\ZZ$, eivät molemmat nollia. Niiden suurin yhteinen
  tekijä $d=\syt(a,b)$ voidaan esittää muodossa
  \[
    d=ua+vb,\qquad u,v\in\ZZ.
  \]
\end{lemma}
\begin{proof}
  Valitaan pienin positiivinen luku $d$ muotoa $ua+vb$.
  Sellainen on olemassa, sillä ainakin toinen luvuista $|a|,|b|$
  on positiivinen ja tätä muotoa. Jakoyhtälössä $a=qd+r$,
  $0\le r<d$, myös $r=a-qd$ on lukujen $a,b$ kokonaislukukertoiminen
  summa. Minimaalisuus pakottaa $r=0$, joten $d\mid a$.
  Samoin $d\mid b$. Jokainen yhteinen tekijä jakaa luvun $ua+vb=d$,
  joten $d$ on suurin positiivinen yhteinen tekijä.
\end{proof}

\begin{lemma}[Eukleideen lemma]
  Jos $p$ on alkuluku ja $p\mid ab$, niin $p\mid a$ tai $p\mid b$.
\end{lemma}
\begin{proof}
  Jos $p\nmid a$, niin $\syt(p,a)=1$, sillä luvun $p$ ainoat
  positiiviset tekijät ovat $1$ ja $p$. Bézout'n identiteetillä
  $up+va=1$ joillakin $u,v\in\ZZ$. Kertomalla luvulla $b$ saadaan
  $b=upb+vab$. Molemmat oikean puolen termit ovat jaollisia luvulla
  $p$, joten $p\mid b$.
\end{proof}
Toistamalla lemma saadaan: jos $p\mid a_1\cdots a_r$, niin
$p\mid a_i$ ainakin yhdellä indeksillä $i$.
Alkulukuoletus on välttämätön: $6\mid2\cdot3$, mutta
$6\nmid2$ ja $6\nmid3$.

\subsubsection{Aritmetiikan peruslause}
\begin{theorem}[Aritmetiikan peruslause]
  Jokainen kokonaisluku $n\ge2$ on alkulukujen tulo.
  Esitys on tekijöiden järjestystä vaille yksikäsitteinen.
  Yhdistämällä samat alkuluvut saadaan \emph{kanoninen alkutekijähajotelma}
  \[
    \boxed{n=p_1^{\alpha_1}\cdots p_k^{\alpha_k},\qquad
    p_1<\cdots<p_k,\quad \alpha_i\ge1,}
  \]
  missä $p_i$ ovat alkulukuja ja $\alpha_i$ kokonaislukuja.
  Sekä alkuluvut että niiden eksponentit määräytyvät yksikäsitteisesti.
\end{theorem}
\begin{proof}
  \emph{Olemassaolo:} käytetään vahvaa induktiota luvun $n\ge2$
  suhteen. Jos $n$ on alkuluku, väite on selvä (erityisesti kun
  $n=2$). Muuten $n=ab$, missä $2\le a,b<n$. Induktio-oletuksen
  mukaan $a$ ja $b$ ovat alkulukujen tuloja, joten myös $n$ on.

  \emph{Yksikäsitteisyys:} olkoon
  $n=p_1\cdots p_r=q_1\cdots q_s$ kaksi alkulukuesitystä,
  joissa sama alkuluku saa toistua. Koska $p_1\mid q_1\cdots q_s$,
  Eukleideen lemma antaa $p_1\mid q_j$ jollakin $j$.
  Alkulukujen vuoksi $p_1=q_j$. Järjestetään tämä tekijä ensimmäiseksi
  ja supistetaan se puolittain. Toistetaan sama jäljelle jääville
  tekijöille. Kummaltakaan puolelta ei voi loppua tekijöitä ennen
  toista, sillä silloin $1$ olisi vähintään yhden alkuluvun tulo.
  Siis $r=s$ ja tekijät ovat samat kertalukuineen.
\end{proof}

\paragraph{Esimerkki: hajotelma ja jaollisuus.}
\[
  360=36\cdot10=2^3\cdot3^2\cdot5.
\]
Sen positiiviset tekijät ovat täsmälleen luvut
$2^a3^b5^c$, missä $0\le a\le3$, $0\le b\le2$ ja $0\le c\le1$.
Nimittäin tekijä ja sen osamäärä muodostavat yhdessä luvun $360$
alkutekijähajotelman; kääntäen näillä eksponenteilla osamäärä on
$2^{3-a}3^{2-b}5^{1-c}\in\ZZ$.
Tekijöitä on siis $4\cdot3\cdot2=24$.
Esimerkiksi $72=2^3\cdot3^2\mid360$, mutta $27=3^3\nmid360$.
Luku $1$ vastaa tyhjää tuloa. Negatiivisen kokonaisluvun
$n\le-2$ hajotelma saadaan liittämällä luvun $|n|$ hajotelman eteen $-1$.

\subsubsection{Alkuluvun testaaminen}
\begin{lemma}[Neliöjuuriraja]
  Jos $n\ge2$ on yhdistetty, sillä on alkutekijä $p\le\sqrt n$.
\end{lemma}
\begin{proof}
  Kirjoitetaan $n=ab$, missä $2\le a\le b$. Tällöin $a^2\le ab=n$,
  joten $a\le\sqrt n$. Luvulla $a$ on alkutekijä $p$, ja
  $p\mid a\mid n$ sekä $p\le a\le\sqrt n$.
\end{proof}
Siis luvun $n\ge2$ alkuluvuksi osoittamiseen riittää tarkistaa,
ettei mikään alkuluku $p\le\sqrt n$ jaa sitä. Raja on suljettu:
myös $p=\sqrt n$ on tarkistettava, kuten $49=7^2$ osoittaa.

\paragraph{Esimerkkejä.}
Koska $9<\sqrt{97}<10$, kokeillaan vain alkulukuja $2,3,5,7$.
Luku $97$ on pariton, sen numeroiden summa on $16$, viimeinen numero
on $7$, ja $97=7\cdot13+6$. Mikään näistä alkuluvuista ei siis jaa
lukua $97$, joten se on alkuluku.
Sen sijaan $221=13\cdot17$ on yhdistetty: alkutekijä $13$ löytyy
ennen rajaa $\sqrt{221}<15$.

\newpage
\subsubsection{Lisäesimerkkejä: rivioperaatiot lukuteoriassa}
Rivioperaatioilla voidaan toteuttaa Eukleideen algoritmi ja samalla
laskea Bézout'n kertoimet. Tämä täydentää luennon alkulukuteoriaa.

\paragraph{Miksi kokonaislukurivioperaatiot toimivat?}
Luvuilla $a,b$ ja $a-qb,b$ on samat yhteiset tekijät, kun $q\in\ZZ$:
\[
  c\mid a,\ c\mid b
  \quad\Longleftrightarrow\quad
  c\mid(a-qb),\ c\mid b.
\]
Suunta takaisin seuraa yhtälöstä $a=(a-qb)+qb$.
Siksi rivien vaihto, merkinvaihto ja operaatio
$R_i\leftarrow R_i-qR_j$ säilyttävät lukujen suurimman yhteisen tekijän.
Jakoyhtälön mukainen vähennys tuottaa aina pienemmän epänegatiivisen
jakojäännöksen. Menettely päättyy nollaan, ja viimeinen positiivinen
jakojäännös on suurin yhteinen tekijä, koska $\syt(d,0)=d$.

\paragraph{Esimerkki 1: suurin yhteinen tekijä ja Bézout'n kertoimet.}
Lasketaan $\syt(252,198)$. Rivissä $(r,u,v)$ pidetään voimassa
$r=252u+198v$; siksi sama operaatio tehdään kaikille kolmelle sarakkeelle.
\[
  \begin{array}{c|rr}
    r&u&v\\\hline
    252&1&0\\198&0&1
  \end{array}
  \xrightarrow{R_1\leftarrow R_1-R_2}
  \begin{array}{c|rr}
    54&1&-1\\198&0&1
  \end{array}
  \xrightarrow{R_2\leftarrow R_2-3R_1}
  \begin{array}{c|rr}
    54&1&-1\\36&-3&4
  \end{array}
\]
\[
  \xrightarrow{R_1\leftarrow R_1-R_2}
  \begin{array}{c|rr}
    18&4&-5\\36&-3&4
  \end{array}
  \xrightarrow{R_2\leftarrow R_2-2R_1}
  \begin{array}{c|rr}
    18&4&-5\\0&-11&14
  \end{array}.
\]
Siis $\syt(252,198)=18$ ja
\[
  \boxed{18=4\cdot252-5\cdot198.}
\]
Sama identiteetti ratkaisee kokonaislukuyhtälön $252x+198y=36$:
kertomalla kahdella saadaan yksi ratkaisu $(x,y)=(8,-10)$.
Yhtälöllä $252x+198y=10$ ei ole kokonaislukuratkaisua, sillä vasen
puoli on aina jaollinen luvulla $18$, mutta oikea puoli ei ole.

\paragraph{Esimerkki 2: käänteisluku ja kongruenssin ratkaiseminen.}
Etsitään luvun $17$ käänteisluku modulo $43$. Nyt rivi $(r,u,v)$
tarkoittaa yhtälöä $r=43u+17v$.
\[
  \begin{array}{c|rr}
    43&1&0\\17&0&1
  \end{array}
  \xrightarrow{R_1\leftarrow R_1-2R_2}
  \begin{array}{c|rr}
    9&1&-2\\17&0&1
  \end{array}
  \xrightarrow{R_2\leftarrow R_2-R_1}
  \begin{array}{c|rr}
    9&1&-2\\8&-1&3
  \end{array}
\]
\[
  \xrightarrow{R_1\leftarrow R_1-R_2}
  \begin{array}{c|rr}
    1&2&-5\\8&-1&3
  \end{array}.
\]
Ensimmäinen rivi antaa $1=2\cdot43-5\cdot17$, joten
$17(-5)\equiv1\pmod{43}$. Käänteisluku on siis $-5\equiv38$.
Näin
\[
  17x\equiv6\pmod{43}
  \quad\Longleftrightarrow\quad
  x\equiv-30\equiv13\pmod{43}.
\]
Tarkistus: $17\cdot13=221=5\cdot43+6$.

\paragraph{Huomaa ero tavalliseen rivireduktioon.}
Suurinta yhteistä tekijää laskettaessa riviä ei saa jakaa mielivaltaisesti:
esimerkiksi parin $(6,10)$ syt on $2$, mutta parin $(3,10)$ syt on $1$.
Kongruenssiyhtälössä rivin saa kertoa kertoimella, jolla on käänteisluku
kyseisen modulon suhteen. Yllä luvulla $17$ jakaminen tarkoitti juuri
kertomista sen käänteisluvulla $38$ modulo $43$.

\paragraph{Kertauksen ydin.}
Jokaisella luvulla $n\ge2$ on alkutekijä, ja alkulukuja on äärettömästi.
Alkutekijähajotelman olemassaolo seuraa vahvasta induktiosta;
yksikäsitteisyys seuraa Eukleideen lemmasta. Alkuluvun testaamisessa
riittävät alkulukujakajat neliöjuureen asti. Kokonaislukurivioperaatiot
antavat syt:n lisäksi konkreettiset Bézout'n kertoimet.

\end{document}
